\documentclass[10pt,a4paper]{article}
\usepackage[margin=16mm,top=17mm,bottom=16mm]{geometry}
\usepackage{fontspec}
\setmainfont{DejaVu Sans}
\usepackage{booktabs,tabularx,array,graphicx,xcolor,hyperref,fancyhdr,enumitem}
\definecolor{navy}{HTML}{1B4965}
\definecolor{red}{HTML}{8E1216}
\definecolor{green}{HTML}{1A5E39}
\definecolor{light}{HTML}{F1F5F7}
\hypersetup{colorlinks=true,linkcolor=navy,urlcolor=navy}
\pagestyle{fancy}\fancyhf{}
\lhead{\small\color{navy}\textbf{MIYAGI - auditoria independente}}
\rhead{\small 13/08/2026}\cfoot{\thepage}
\setlength{\parindent}{0pt}\setlength{\parskip}{5pt}
\setlist[itemize]{leftmargin=5mm,itemsep=2pt,topsep=3pt}
\setlist[enumerate]{leftmargin=6mm,itemsep=2pt,topsep=3pt}
\newcommand{\critical}[1]{\colorbox{red}{\textcolor{white}{\strut\textbf{#1}}}}
\newcommand{\high}[1]{\colorbox{navy}{\textcolor{white}{\strut\textbf{#1}}}}
\newcommand{\ok}[1]{\colorbox{green}{\textcolor{white}{\strut\textbf{#1}}}}
\newenvironment{callout}{\begin{center}\begin{minipage}{0.95\linewidth}\color{navy}\hrule height 1.5pt\vspace{2mm}\color{black}}{\vspace{1mm}\color{navy}\hrule height 0.5pt\end{minipage}\end{center}}

\begin{document}

{\Huge\color{navy}\textbf{MIYAGI}}\\[3mm]
{\Large Auditoria crítica e independente}\\[2mm]
{\small Estratégia de time-series momentum multiativo - Desafio Quant AI 2026}

\begin{callout}
\textbf{Veredito.} A estratégia continua positiva e passa os seis critérios de
robustez declarados. Porém, o Sharpe auditado é \textbf{0,47}, não 0,50, e a
evidência estatística é marginal. O histórico não pode ser chamado integralmente
fora da amostra, e o painel mistura instrumentos com economias incompatíveis.
Este é um proxy acadêmico de pesquisa, não ainda um retorno implementável.
\end{callout}

\section*{Resumo quantitativo}
\begin{tabularx}{\textwidth}{Xrrrr}
\toprule
Versão & CAGR & Vol. & Sharpe & Max DD\\
\midrule
Relatório anterior & 15,9\% & 10,5\% & 0,50 & -23,0\%\\
\textbf{Motor auditado} & \textbf{15,6\%} & \textbf{10,5\%} & \textbf{0,47} & \textbf{-24,1\%}\\
Seleção anual defasada, 2016-2026 & 12,9\% & 10,6\% & 0,34 & -25,0\%\\
\bottomrule
\end{tabularx}

\textbf{O que mudou no motor:} retorno calculado com
\texttt{fill\_method=None}; painel corrigido usado em todos os scripts; pesos
efetivos derivam entre rebalanceamentos; contribuição usa o peso anterior ao
retorno; lacunas e exposição sem retorno são registradas.

\section*{Achados prioritários}
\begin{enumerate}
\item \critical{CRÍTICO} \textbf{Linhagem inconsistente.} Robustez e figuras
usavam \texttt{pool\_expandido.csv}, enquanto o resultado final usava
\texttt{pool\_carrego.csv}. A carteira publicada mostrava TRY comprada; o motor
corrigido a coloca vendida em -0,614.
\item \critical{CRÍTICO} \textbf{Seleção com informação futura.} Elegibilidade,
correlações e medoides foram calculados até 2026 e aplicados desde 2005. A regra
de 15 anos nem poderia ser executada em 2005 sobre um pool iniciado em 2000.
\item \critical{CRÍTICO} \textbf{Dados ausentes dependiam da versão do pandas.}
O padrão antigo de \texttt{pct\_change()} preenchia buracos e criava saltos de
reabertura. A opção explícita sem preenchimento reproduzia o 0,50 documentado,
mas ainda deixava posições ativas durante lacunas.
\item \high{ALTO} \textbf{Economia dos instrumentos.} O painel soma CDI a ETFs,
índices de preço, futuros contínuos e FX. Essa soma só é correta se todas as
pernas arriscadas forem retornos excedentes financiados, o que não ocorre.
\item \high{ALTO} \textbf{Inferência frágil.} Autocorrelação, caudas e o caminho
de pesquisa tornam t=2,19 insuficiente para uma alegação forte de significância.
\end{enumerate}

\newpage
\section*{1. Modelagem e execução}

\subsection*{1.1 Rebalanceamento mensal versus diário implícito}
O código anterior aplicava o mesmo vetor de pesos todos os dias do mês. Isso
equivale a restaurar diariamente os pesos-alvo, sem cobrar as negociações dessa
restauração. O motor auditado mantém quantidades/notionais entre datas mensais e
deixa os pesos derivarem com preços e patrimônio. A mudança reduz o Sharpe de
0,50 para 0,47 e aumenta o drawdown de -23,0\% para -24,1\%.

\subsection*{1.2 Atribuição temporal}
O relatório atribuía o retorno do dia de rebalanceamento à carteira nova. A
estratégia declara que decide no fechamento e só colhe retornos a partir do dia
seguinte. O novo log \texttt{pesos\_diarios} registra a exposição efetiva antes
de cada retorno. TRY continua sendo a maior contribuição, agora cerca de
+1,16\% ao ano, mas a última posição correta é vendida.

\subsection*{1.3 Lacunas}
Foram identificadas \textbf{60 lacunas internas maiores que cinco dias}. Em
\textbf{330 de 5.386 dias} de carregamento (6,1\%) havia posição ativa e ao menos
um retorno ausente. A exposição absoluta média nesses dias foi 6,5\%; o máximo,
44,2\%. O motor mantém P\&L zero onde não há retorno e registra a exposição, mas
isso é uma convenção de diagnóstico, não uma cotação observada.

\begin{callout}
\textbf{Decisão correta:} não interpolar preços nem inventar retornos. Para uma
versão implementável, cada lacuna precisa ser resolvida na fonte ou levar a uma
regra de saída/tradabilidade definida antes do teste.
\end{callout}

\subsection*{1.4 Instrumentos}
\begin{tabularx}{\textwidth}{l r X}
\toprule
Tipo no universo & N & Problema econômico\\
\midrule
Futuros contínuos Yahoo & 15 & roll e custo de rolagem não identificados\\
ETFs adjusted close & 11 & retorno total financiado, não excesso de futuro\\
FX com carry proxy & 6 & taxas proxy; sem vintage e lag de publicação\\
FX spot & 4 & carry ausente\\
Índices de preço & 4 & dividendos e conversão cambial ausentes\\
\bottomrule
\end{tabularx}

Adicionar CDI integralmente é defensável para uma carteira de futuros com
collateral, desde que cada série arriscada seja um retorno excedente de contrato.
Não é defensável de forma uniforme para o painel atual. Também faltam roll,
borrow de posições vendidas, margem, haircuts e custos específicos.

\newpage
\section*{2. Viés de seleção}

\subsection*{2.1 O vazamento}
O funil é cego a retorno médio, mas não ao futuro. Um ativo só entra se, visto em
2026, tiver 15 anos de histórico; a correlação e o representante do cluster são
escolhidos sobre toda a amostra. Saber antecipadamente quais ativos sobreviverão
e como se correlacionarão é informação futura.

\subsection*{2.2 Teste anual defasado}
Foi implementada uma seleção no último dia disponível de cada ano, usada apenas
no ano seguinte. Mantém-se a regra declarada de 15 anos, sem encurtá-la para
obter um resultado melhor. Por isso o primeiro corte possível é dezembro de
2015 e o primeiro retorno é de 2016.

\begin{tabularx}{\textwidth}{Xrrrr}
\toprule
Período & CAGR & Vol. & Sharpe & Max DD\\
\midrule
2016-2026, seleção anual defasada & 12,9\% & 10,6\% & 0,34 & -25,0\%\\
\bottomrule
\end{tabularx}

O número é positivo, mas não é comparável diretamente ao Sharpe 0,47 de
2005-2026. Ele responde a uma pergunta mais restrita: quanto resta quando o
funil explícito só usa o passado, no período em que o requisito é possível.

\begin{callout}
\textbf{Rótulo obrigatório: pseudo-point-in-time.} Os CSVs atuais foram
baixados retrospectivamente. Sem snapshots/vintages do Yahoo e ALFRED não se
prova que o provedor mostrava exatamente esses valores e esse conjunto de ativos
em cada data histórica.
\end{callout}

\subsection*{2.3 Pré-registro}
O commit \texttt{88392ef} registrou Sharpe 0,65 para o walk-forward de 40 ativos.
O realizado daquela etapa foi 0,48, fora da banda de ±0,10. Logo, a previsão
original falhou. Recalcular depois a fórmula com Sharpe 0,36 e 11,2 apostas para
obter 0,455, comparando com a amostra completa, é uma checagem pós-hoc de
consistência, não confirmação independente.

\section*{3. Inferência estatística}
\begin{tabularx}{\textwidth}{lrrr}
\toprule
Estimador & Lag & t & p bilateral\\
\midrule
Diário IID & 0 & 2,189 & 0,0286\\
Diário Newey-West & 21 & 1,908 & 0,0564\\
Diário Newey-West & 252 & 1,936 & 0,0529\\
Mensal IID & 0 & 1,988 & 0,0468\\
Mensal Newey-West & 24 & 1,913 & 0,0557\\
\bottomrule
\end{tabularx}

\newpage
\section*{4. Múltiplos testes e incerteza}

O p bilateral IID é 0,0286. Bonferroni para apenas quatro combinações dá
\textbf{0,114}; no teste unilateral, 0,057. Quatro subestima o caminho total,
que inclui universos, estimadores, horizontes, janelas, alavancagem, expansão e
mudanças de modelo. Portanto nem esse ajuste deve ser apresentado como
conservador completo.

\begin{tabularx}{\textwidth}{lrrr}
\toprule
Bloco bootstrap & IC 2,5\% & Mediana & IC 97,5\%\\
\midrule
21 dias & -0,037 & 0,466 & 0,947\\
63 dias & 0,025 & 0,470 & 0,960\\
252 dias & 0,002 & 0,486 & 0,960\\
\bottomrule
\end{tabularx}

Os três blocos foram declarados e reportados conjuntamente; nenhum foi
selecionado por favorecer a estratégia. O limite inferior muda de sinal, o que
reforça a leitura de evidência marginal. Um passo adicional adequado é White
Reality Check/SPA ou Deflated Sharpe Ratio sobre o conjunto completo de
estratégias realmente tentadas, cuja contagem precisa ser reconstruída do log.

\section*{5. Robustez após as correções}
\begin{tabularx}{\textwidth}{l X r}
\toprule
Teste & Resultado auditado & Veredito\\
\midrule
Custos & 0,2\%/trade: CAGR 14,3\% & \ok{PASSA}\\
Janela de vol & Sharpe mínimo 0,40 & \ok{PASSA}\\
Subperíodos & supera CDI em 3 de 4 & \ok{PASSA}\\
Alavancagem & teto 2x: CAGR 14,1\% & \ok{PASSA}\\
Jackknife por classe & pior Sharpe 0,41, sem câmbio & \ok{PASSA}\\
Horizonte & Sharpe positivo em 4 de 4 & \ok{PASSA}\\
\bottomrule
\end{tabularx}

Isso permite dizer: \textit{o resultado não depende de uma única escolha entre
as seis variações testadas}. Não permite dizer: \textit{o resultado está livre
de erro de modelagem, seleção ou múltiplos testes}. O erro de carry original já
havia passado por testes de robustez, mostrando por que robustez paramétrica e
validade econômica são perguntas diferentes.

\section*{6. Controle de risco}
O teto de 3x é atingido em \textbf{48,3\%} dos rebalanceamentos, não cerca de
23\%. Logo, ele é parte ativa do perfil de retorno e risco, embora o teste com
2x continue superando o CDI. A volatilidade agregada de 10,5\% próxima ao alvo
de 10\% não valida sozinha a previsão mensal de risco nem a matriz 40x40
estimada com apenas 60 dias.

\newpage
\section*{7. O que pode ser afirmado}

{\small\begin{tabularx}{\textwidth}{>{\raggedright\arraybackslash}X >{\raggedright\arraybackslash}X}
\toprule
Afirmação sustentada & Afirmação que deve ser removida\\
\midrule
A estratégia tem retorno excedente positivo na amostra. & ``Sharpe 0,50 é
estatisticamente robusto.''\\
Os seis critérios declarados sobrevivem ao motor auditado. & ``2005-2026 é todo
fora da amostra.''\\
O funil anual causal mantém Sharpe positivo em 2016-2026. & ``A previsão de
diversificação foi confirmada.''\\
TRY continua relevante após a correção de carry. & ``A queda da duração das
tendências causou o bloco ruim.''\\
O painel serve a uma prova acadêmica transparente. & ``O backtest representa
execução real de futuros.''\\
\bottomrule
\end{tabularx}}

\section*{8. Ordem de correção recomendada}
\begin{enumerate}
\item Usar uma única fonte oficial em código, gráficos, robustez e carteira.
\item Fixar versões, retorno sem preenchimento e auditoria de lacunas.
\item Manter timing causal e deriva de pesos entre rebalanceamentos.
\item Rotular o funil completo como in-sample e reportar o teste anual 2016+.
\item Reescrever significância e pré-registro com HAC e múltiplos testes.
\item Reconstruir instrumentos: contratos e rolls de futuros; forwards de FX;
índices total return ou futuros de ações; funding e collateral consistentes.
\end{enumerate}

\section*{Artefatos reproduzíveis}
\begin{itemize}
\item \texttt{dados\_miyagi.py}: fonte única e regras de calendário.
\item \texttt{auditoria\_dados.py}: lacunas, exposição sem retorno e tipos.
\item \texttt{auditoria\_estatistica.py}: IID, HAC, Bonferroni e bootstrap.
\item \texttt{auditoria\_selecao.py}: funil anual defasado.
\item \texttt{robustez.py --universo 40}: seis testes no painel oficial.
\item \texttt{tests/test\_backtest.py}: testes contra preenchimento implícito,
sinal sem cotação e divergência de painel.
\end{itemize}

\begin{callout}
\textbf{Conclusão final.} O projeto tem mérito por registrar erros e manter
parâmetros mesmo quando alternativas parecem melhores. A contribuição mais
forte para o desafio não é vender um Sharpe de 0,50 como fato; é mostrar uma
pesquisa que encontrou onde o próprio número falhava, corrigiu o motor e
preservou as limitações sem fabricar dados.
\end{callout}

\vfill
{\small\color{navy}\textbf{Fontes metodológicas principais}}\\
{\footnotesize Moskowitz, Ooi e Pedersen (2012), \textit{Time Series Momentum};
Hurst, Ooi e Pedersen (2017), \textit{A Century of Evidence on Trend-Following};
White (2000), \textit{A Reality Check for Data Snooping}; Bailey e López de Prado
(2014), \textit{The Deflated Sharpe Ratio}. Resultados numéricos desta auditoria
foram produzidos exclusivamente pelos CSVs versionados no repositório.}

\end{document}
